% =====================================================================
%  Stage 2, раздел 2.
%  Теоретическая модель данных и pipeline обработки.
% =====================================================================

\section{Теоретическая модель данных и pipeline обработки}

\subsection{Схема признаков UNSW-NB15}

UNSW-NB15~\cite{moustafa2015unsw, moustafa2016eval, unswportal} в составе
\emph{cleaned}-распределения содержит 49 признаков, разделяемых на пять
групп. Назначение каждой группы и её роль в детекции атак приведены в
\tabref{tab:unsw-feature-groups}.

\begin{table}[H]
\centering
\caption{Группы признаков UNSW-NB15 и их назначение}
\label{tab:unsw-feature-groups}
\renewcommand{\arraystretch}{1.25}
\begin{tabularx}{\linewidth}{@{}lXX@{}}
\toprule
\textbf{Группа} & \textbf{Состав} & \textbf{Роль в детекции} \\
\midrule
Базовые потоковые &
\emph{dur, proto, service, state, sbytes, dbytes, sttl, dttl, sloss, dloss}
&
непосредственное описание flow: длительность, протокол, состояние сессии,
байты и пакеты в обе стороны \\
\addlinespace
Контентные &
\emph{sload, dload, swin, dwin, smean, dmean, stcpb, dtcpb, trans\_depth,
response\_body\_len} &
скорости, окна TCP, базовые транзакционные характеристики \\
\addlinespace
Временные &
\emph{sjit, djit, sintpkt, dintpkt, tcprtt, synack, ackdat} &
межпакетные интервалы, джиттеры, RTT, timing-сигнатуры \\
\addlinespace
Дополнительные сгенерированные &
\emph{ct\_state\_ttl, ct\_flw\_http\_mthd, is\_ftp\_login,
ct\_ftp\_cmd} &
агрегаты по протокольному поведению \\
\addlinespace
Контекстные \emph{ct\_*}&
\emph{ct\_srv\_src, ct\_srv\_dst, ct\_dst\_ltm, ct\_src\_ltm,
ct\_src\_dport\_ltm, ct\_dst\_sport\_ltm, ct\_dst\_src\_ltm} &
агрегаты по последним 100 соединениям; косвенно отражают режим
сканирования и фан-аут \\
\bottomrule
\end{tabularx}
\end{table}

Метки представлены двумя полями: \emph{label} \(\in \{0, 1\}\) (бинарная) и
\emph{attack\_cat} \(\in \mathcal{C} \cup \{\varnothing\}\) (мультиклассовая,
с пустым значением для нормального трафика).

\subsection{Схема признаков CICIDS2017 и подмножество для cross-evaluation}

CICIDS2017~\cite{sharafaldin2018toward} формируется утилитой CICFlowMeter и
содержит более 80 flow-признаков. Часть из них функционально эквивалентна
полям UNSW-NB15 (длительность, протокол, число пакетов и байт в обе
стороны, межпакетные интервалы, скорости передачи). Подмножество общих
признаков, принимаемое в работе, представлено в
\tabref{tab:common-features}.

\begin{table}[H]
\centering
\caption{Подмножество признаков, общих для UNSW-NB15 и CICIDS2017,
принятое для cross-evaluation}
\label{tab:common-features}
\renewcommand{\arraystretch}{1.25}
\begin{tabularx}{\linewidth}{@{}lll@{}}
\toprule
\textbf{Семантика} & \textbf{UNSW-NB15} & \textbf{CICIDS2017 (CICFlowMeter)} \\
\midrule
Длительность потока       & \emph{dur}      & \emph{Flow Duration} \\
Протокол транспорта        & \emph{proto}    & \emph{Protocol} \\
Прикладной сервис          & \emph{service}  & \emph{Destination Port}+\emph{Service heuristics} \\
Байты от источника         & \emph{sbytes}   & \emph{Total Length of Fwd Packets} \\
Байты к источнику          & \emph{dbytes}   & \emph{Total Length of Bwd Packets} \\
Пакеты от источника        & \emph{spkts}    & \emph{Total Fwd Packets} \\
Пакеты к источнику         & \emph{dpkts}    & \emph{Total Backward Packets} \\
TTL источника              & \emph{sttl}     & ---  (косвенно через \emph{Init\_Win\_bytes\_forward})\\
Скорость источника         & \emph{sload}    & \emph{Fwd Packets/s} \\
Скорость получателя        & \emph{dload}    & \emph{Bwd Packets/s} \\
Межпакетный интервал, src  & \emph{sintpkt}  & \emph{Fwd IAT Mean} \\
Межпакетный интервал, dst  & \emph{dintpkt}  & \emph{Bwd IAT Mean} \\
Джиттер, src               & \emph{sjit}     & \emph{Fwd IAT Std} \\
Джиттер, dst               & \emph{djit}     & \emph{Bwd IAT Std} \\
SYN-ACK / ACK-DAT timings  & \emph{synack, ackdat} & \emph{Flow IAT}-агрегаты \\
\bottomrule
\end{tabularx}
\end{table}

Признаки, не имеющие прямого аналога (TTL, ряд \emph{ct\_*}-агрегатов
UNSW-NB15, специфические признаки CICFlowMeter), исключаются из набора
cross-evaluation. Эта позиция фиксирована и применяется в эксперименте E4
протокола (раздел~3 Этапа~2).

\subsection{Pipeline обработки данных}

Полный pipeline обработки данных принимает на вход поток flow-записей и
формирует на выходе скоринг \(s(x)\) и решение \(\hat{y}\):

\[
\boxed{\;
\text{IPFIX/CSV} \to
\text{Schema-Valid} \to
\text{Cleanup} \to
\text{Encoding} \to
\text{Scaling} \to
\text{Windowing} \to
\text{Model} \to
\text{Threshold} \to
\text{Alert}
\;}
\]

Стадия каждого блока формализуется ниже.

\subsection{Очистка данных}

В UNSW-NB15 документированы~\cite{moustafa2015unsw} следующие виды
``шума'': пропуски в полях \emph{ct\_flw\_http\_mthd, is\_ftp\_login,
ct\_ftp\_cmd}, неконсистентные строковые значения \emph{service} и
\emph{state}, отдельные крайние выбросы в \emph{dur, sbytes, dbytes}.
Применяемая очистка:

\begin{enumerate}
    \item Замена нечисловых пропусков константой \texttt{"-"} (отсутствие
    значения). Для числовых пропусков --- замена нулём, что согласуется с
    интерпретацией признаков (отсутствие активности).
    \item Унификация регистра и пробелов в \emph{service}, \emph{state},
    \emph{proto}.
    \item Удаление дубликатов по полному вектору признаков на стороне
    обучающей выборки (исключение перекрёстной утечки).
    \item Идентификация и обрезка выбросов в крайних \(0{,}1\%\) хвостов
    распределений для \emph{dur, sbytes, dbytes, sload, dload}; границы
    обрезки фиксируются по обучающей выборке и применяются к валидационной
    и тестовой без пересчёта.
\end{enumerate}

Замечание: вся статистика для очистки и масштабирования вычисляется
\emph{только} на обучающей выборке, что исключает test-time leakage.

\subsection{Кодирование категориальных признаков}

\subsubsection{One-hot}

Для признака \(c\) с конечным алфавитом \(\mathcal{A}_c =
\{a_1, \dots, a_K\}\) one-hot-вектор имеет вид
\(e(c) = (\mathbb{I}[c=a_1], \dots, \mathbb{I}[c=a_K])^\top \in \{0,1\}^K\).
Применимо для признаков с малой кардинальностью --- \emph{proto, state}.

\subsubsection{Frequency-encoding и target-encoding}

Для признака с большой кардинальностью \emph{service}
(\(\sim 13\) уровней) применяется frequency-encoding:
\begin{equation}
e_{\text{freq}}(c) = \frac{|\{i : x_i^{(c)} = c\}|}{n_{\text{train}}},
\label{eq:freq-enc}
\end{equation}
а как опциональная альтернатива --- target-encoding со сглаживанием:
\begin{equation}
e_{\text{tgt}}(c) =
\frac{\sum_i \mathbb{I}[x_i^{(c)}=c]\,y_i + \alpha\,\bar{y}}
     {\sum_i \mathbb{I}[x_i^{(c)}=c] + \alpha},
\label{eq:tgt-enc}
\end{equation}
где \(\bar{y}\) --- средняя метка класса в обучающей выборке, \(\alpha\) ---
параметр сглаживания. Target-encoding применяется с осторожностью и только
при кросс-валидационной схеме во избежание утечки.

\subsubsection{Embedding categorical}

В DL-моделях категориальные признаки могут кодироваться обучаемыми
embeddings \(E_c \in \mathbb{R}^{K \times d_e}\), где \(d_e\) --- размерность
embedding. Это совместимо с архитектурой CNN-LSTM и, по обзорам
NIDS-DL~\cite{ferrag2020dlcyber, ahmad2021nids}, даёт прирост качества
относительно one-hot для признаков большой кардинальности.

\subsection{Преобразования числовых признаков}

\subsubsection{Логарифмическое преобразование}

Признаки с положительной асимметрией и тяжёлыми хвостами (\emph{dur, sbytes,
dbytes, sload, dload, sintpkt, dintpkt}) подвергаются преобразованию
\begin{equation}
\tilde{x} = \log(1 + x).
\label{eq:log1p}
\end{equation}
\(\log(1+x)\) корректно обрабатывает нулевые значения и эквивалентна
гомеоморфизму, переводящему распределение в близкое к нормальному.

\subsubsection{Robust-scaling}

После логарифмирования признаки масштабируются робастно:
\begin{equation}
\tilde{x}' = \frac{\tilde{x} - \operatorname{med}(\tilde{x})}
                  {\operatorname{IQR}(\tilde{x})},
\label{eq:robust-scaler}
\end{equation}
где \(\operatorname{med}\) --- медиана, \(\operatorname{IQR}\) ---
интерквартильный размах, оцениваемые \emph{на обучающей выборке}.
Робастные оценки устойчивы к выбросам, что критично для NIDS, где
аномалии формируют ``хвосты'' распределения.

\subsubsection{Стандартизация}

Для baseline-моделей (LR, MLP) дополнительно используется стандартизация
\(z = (x - \mu)/\sigma\). Для DL-моделей эта операция эквивалентна robust-scaling
с заменой \(\operatorname{med} \to \mu\), \(\operatorname{IQR} \to \sigma\).

\subsection{Формирование скользящих окон}

Для рекуррентных и гибридных моделей (CNN-LSTM, BiLSTM, Transformer)
исходный поток flow-записей преобразуется в скользящие окна фиксированной
длины \(W\) с шагом \(S\):
\[
\mathcal{W}_t = \bigl(x_{t-W+1}, x_{t-W+2}, \dots, x_t\bigr) \in \mathbb{R}^{W \times d},
\]
где \(d\) --- размерность признакового вектора после препроцессинга и
кодирования.

\subsubsection{Выбор параметров \texorpdfstring{$W$ и $S$}{W и S}}

Параметр \(W\) определяет компромисс между разрешающей способностью по
времени (короткие \(W\) лучше реагируют на мгновенные всплески) и
полнотой контекста (длинные \(W\) лучше учитывают slow-атаки). По данным
обзоров~\cite{vinayakumar2019dl, ferrag2020dlcyber, ahmad2021nids}
и эталонной работы Kim и соавт.~\cite{kim2016lstm} распространённый диапазон
для NIDS-окон в flow-представлении --- \(W \in \{8, 16, 32, 64\}\).
В настоящей работе принят базовый \(W = 32\) с проверкой устойчивости при
\(W \in \{16, 32, 64\}\) в рамках Эксперимента~1 Этапа~4.

Шаг \(S\) для основного режима принят равным \(S = W/4\) (75\% перекрытие).
Это согласуется с практикой потокового анализа в NTA: малое \(S\) уменьшает
TTD, большое --- сокращает вычислительные затраты.

\subsubsection{Целевая разметка окна}

Для бинарной задачи метка окна формируется как
\begin{equation}
y_{\mathcal{W}_t} =
\mathbb{I}\!\left[\sum_{i=t-W+1}^{t} y_i > 0\right],
\label{eq:window-label}
\end{equation}
то есть окно считается атакующим при наличии хотя бы одной атакующей
записи. Этот вариант принят как основной (\emph{any-positive}); как
альтернатива в Эксперименте~1 рассматривается \emph{majority-positive}
(порог по доле атакующих записей).

\subsubsection{Запрет утечки и временной split}

Train/val/test-разделение выполняется при сохранении временного порядка:
обучающая выборка предшествует валидационной, валидационная ---
тестовой. При этом запрещено формирование окна, содержащего записи из
разных сплитов, --- в противном случае возникает утечка
``из будущего''~\cite{pendlebury2019tesseract}.

\subsection{Управление дисбалансом классов}

UNSW-NB15 в \emph{cleaned}-распределении имеет умеренный дисбаланс
классов с доминированием \emph{Generic}. Применяемые меры:

\subsubsection{Re-weighting}

Веса классов \(w_c\) обратно пропорциональны их частотам:
\begin{equation}
w_c = \frac{n}{|\mathcal{C}|\,n_c},
\label{eq:reweight}
\end{equation}
где \(n_c\) --- число примеров класса \(c\) в обучающей выборке. Веса
включаются в функцию потерь.

\subsubsection{Focal loss}

Для усиления внимания к ``тяжёлым'' примерам применяется фокальная функция
потерь~\cite{lin2017focal}:
\begin{equation}
\mathcal{L}_{\text{focal}} =
- \frac{1}{n}\sum_{i=1}^{n}
  \alpha_{y_i}\,(1-p_{y_i})^{\gamma}\,\log p_{y_i},
\label{eq:focal}
\end{equation}
где \(p_{y_i}\) --- предсказанная вероятность истинного класса,
\(\alpha\), \(\gamma\) --- параметры фокусировки. Стандартное значение
\(\gamma = 2{,}0\)~\cite{lin2017focal}.

\subsubsection{SMOTE}

Для классических baseline (LR, RF, XGBoost) применяется
SMOTE~\cite{chawla2002smote}: для каждой точки миноритарного класса
\(x_i\) и её \(k\)-ближайшего соседа того же класса \(x_j\) формируется
синтетическая точка
\begin{equation}
x_{\text{new}} = x_i + \lambda\,(x_j - x_i),\quad \lambda \sim U(0, 1).
\label{eq:smote}
\end{equation}
SMOTE не применяется к окончательной DL-обучающей выборке, так как при
последовательной структуре окон вносит синтетические разрывы во временную
ось.

\subsection{Контроль дрейфа распределений}

В соответствии с разделом~3 Этапа~1 в pipeline предусматривается отдельный
блок drift-monitoring, реализующий метрики:
\begin{itemize}
    \item \(\PSI\) над выбранными признаками (\eqnref{eq:psi});
    \item \(\KL\)-расхождение над дискретизованными распределениями;
    \item \(\MMD\) над непрерывными признаками
    (\eqnref{eq:mmd}).
\end{itemize}

Эталонное окно --- начальное окно нормального трафика в обучающей выборке.
Текущее окно --- скользящее окно того же размера в потоке инференса. При
превышении пороговых значений хотя бы по одной из метрик активируется
расширенный режим логирования и (опционально) переобучение модели.
Конкретные пороги \(\tau_{\PSI} = 0{,}25\), \(\tau_{\MMD}\) (определяется
по 95-перцентилю распределения MMD на bootstrap-выборках обучающего
блока) фиксированы и применяются в Эксперименте~5 Этапа~4.

\subsection{Спецификация форматов входов и выходов}

В \tabref{tab:io-spec} приведены спецификации входных и выходных форматов
системы.

\begin{table}[H]
\centering
\caption{Спецификации входных и выходных данных системы}
\label{tab:io-spec}
\renewcommand{\arraystretch}{1.25}
\begin{tabularx}{\linewidth}{@{}llX@{}}
\toprule
\textbf{Назначение} & \textbf{Формат} & \textbf{Содержание} \\
\midrule
Сырой поток flow-записей &
CSV / JSON Lines &
Поля UNSW-NB15 (49) с типами по \emph{schema.json}; разделитель ``,''; UTF-8 \\
\addlinespace
Препроцессированный поток &
NPZ / Parquet &
Числовая матрица \(\mathbb{R}^{n \times d}\) после кодирования и масштабирования \\
\addlinespace
Окна последовательности &
NPZ / Parquet &
Тензор \(\mathbb{R}^{m \times W \times d}\) с метками \(y \in \{0,1\}^m\) \\
\addlinespace
Скоринг детектора &
JSON Lines &
\texttt{\{ts, score, decision, attack\_cat\_pred, prob, window\_id\}} \\
\addlinespace
Алерт &
JSON Lines / SIEM-event &
\texttt{\{alert\_id, ts, severity, score, profile, sample\_features\}} \\
\addlinespace
Drift-метрика &
JSON Lines &
\texttt{\{ts, psi, kl, mmd, ref\_window, cur\_window\}} \\
\bottomrule
\end{tabularx}
\end{table}

\subsection{Воспроизводимость и контроль конфигурации}

Все шаги pipeline параметризуются конфигурационными YAML-файлами в
директории \texttt{configs/}:

\begin{itemize}
    \item \texttt{configs/data/unsw\_nb15.yaml} --- пути, схема, параметры
    очистки и кодирования;
    \item \texttt{configs/data/cicids2017.yaml} --- то же для CICIDS2017;
    \item \texttt{configs/preprocess/default.yaml} --- параметры
    масштабирования, encoding, windowing;
    \item \texttt{configs/models/cnn\_lstm.yaml} --- архитектура и
    гиперпараметры основной модели;
    \item \texttt{configs/train/default.yaml} --- параметры обучения,
    seed, логи;
    \item \texttt{configs/eval/default.yaml} --- метрики, пороги, схема
    бутстрап-CI.
\end{itemize}

Каждый запуск экспериментов фиксирует:
\begin{enumerate}
    \item версию кода (commit hash);
    \item набор seed (\(\geq 5\));
    \item гиперпараметры обучения и препроцессинга;
    \item версии зависимостей (\texttt{pip freeze});
    \item полное содержимое использованных конфигов в виде attached-артефактов.
\end{enumerate}

\subsection{Промежуточные выводы по разделу}

\begin{enumerate}
    \item Полностью описана схема признаков UNSW-NB15 (5 групп, 49
    признаков) и зафиксировано подмножество, общее с CICIDS2017
    (12--14 признаков), пригодное для протокола cross-evaluation.
    \item Сформирован полный pipeline обработки данных:
    schema-validation $\to$ cleanup $\to$ encoding (one-hot, frequency,
    target, embedding) $\to$ scaling (log, robust) $\to$ windowing
    (\(W=32\), \(S=W/4\)) $\to$ model $\to$ threshold $\to$ alert.
    Каждая стадия снабжена математическим описанием.
    \item Зафиксированы меры управления дисбалансом классов: re-weighting,
    focal loss с параметрами \(\alpha\), \(\gamma\), и SMOTE для классических
    baseline. Введён запрет временной утечки при формировании окон.
    \item Описан блок контроля дрейфа на основе метрик PSI, KL, MMD с
    конкретными пороговыми значениями для эксперимента E5; определены
    эталонное и текущее окна.
    \item Зафиксированы спецификации входных и выходных форматов на всех
    стадиях pipeline и принципы конфигурационного контроля воспроизводимости.
\end{enumerate}
